\documentclass[14pt]{extarticle}
\usepackage{preamble}
\geometry{letterpaper, landscape, margin=0.5in}
\usepackage{qrcode}
\renewcommand{\answer}[1]{}
\setlength{\parskip}{4pt}

\begin{document}
\pagestyle{empty}

\daybanner{Unit 4 \textperiodcentered\ Topic 4.1 \textperiodcentered\ TODAY'S PROBLEM}{Two Production Lines, Two Sampling Plans}

\noindent\begin{minipage}[t]{0.57\linewidth}
\vspace{0pt}
Suppose Lines A and B each contain 1{,}000 cycle times with population mean 40 seconds and standard deviation 12 seconds. Line A is approximately normal and an SRS of $n_A=9$ is selected; Line B is strongly right-skewed and an SRS of $n_B=36$ is selected. An auditor wants $P(\bar X>46)$ for each.\par\smallskip \textbf{Serena:} ``Only B can use a normal model because $36\geq30$; its probability is $0.00135$. A is too small.''\par \textbf{Mateo:} ``Only A can use a normal model because its population is normal; its probability is $0.0668$. B is too skewed.''\par
\end{minipage}\hfill
\begin{minipage}[t]{0.40\linewidth}
\vspace{0pt}
\begin{center}
\textbf{Sampling plans}\par\smallskip
\begin{tabular}{|l|c|c|c|c|c|}
\hline
\textbf{Line} & \textbf{Shape} & \textbf{$N$} & \textbf{$n$} & \textbf{$\mu$} & \textbf{$\sigma$} \\ \hline
\textbf{A} & Normal & 1000 & 9 & 40 & 12 \\ \hline
\textbf{B} & Right-skewed & 1000 & 36 & 40 & 12 \\ \hline
\end{tabular}
\end{center}
\end{minipage}

\begin{firsttakebox}{FIRST TAKE (5 min, on your sheet)}
Which plan seems more likely to give averages close to 40 seconds? Give one reason from the sample sizes or shapes.
\end{firsttakebox}

\begin{description}[leftmargin=1.15in, labelwidth=1in, itemsep=2pt, topsep=2pt]
  \item[\dokbadge{1}~(a)] State the mean of each sampling distribution and the 10\% condition for using $\sigma_{\bar X}=\sigma/\sqrt n$.
  \item[\dokbadge{2}~(b)] Verify the 10 percent condition, compute both sampling-distribution standard deviations, and calculate each normal tail probability when justified.
  \item[\dokbadge{3}~(c)] In 3--4 sentences, explain what each student gets right and wrong. Use the two normality pathways to decide which probabilities are supported. Connect the different tail probabilities to the sampling SDs in (b).
\end{description}

\vfill
\noindent\begin{minipage}{0.84\linewidth}
\textbf{Video + follow-along:} \href{https://robjohncolson.github.io/apstats-live-worksheet/u5_lesson7_live.html}{\texttt{\detokenize{u5_lesson7_live.html}}} (scan the code)\par
\textbf{Finish (a)--(c) after the video. Turn in your sheet.}
\end{minipage}\hfill
\qrcode[height=0.75in]{https://robjohncolson.github.io/apstats-live-worksheet/u5_lesson7_live.html}

\end{document}
